\begin{abstract}
The layout algebra behind NVIDIA's CuTe library describes how tensor tiles map
onto threads and memory. A layout is a shape with a stride for each entry, so a
coordinate's address is the sum of its indices times their strides. As
formalized by Shah and by Cecka, its central operations (e.g., composition,
complement, division, product) are guarded by at least six interacting
admissibility conditions. We show that this complexity is an artifact of a
representational choice, not of the domain.

We present a layout algebra in which every layout names the space its inputs
and outputs live in: logical, physical, or thread-value indices. A layout into
logical or thread-value space must be a bijection from its $N$ coordinates onto
$[0,N)$. A layout into physical space must be alias-free: distinct coordinates
land on distinct addresses, except under broadcast, where many threads read one
value on purpose and the layout must mark where. Nothing maps out of a physical
index.

Composing $f$ with $g$ requires only that $f$'s range be exactly $g$'s
flattened coordinate set. The index arithmetic can always be emitted naively:
compute $f$, split the result by division and remainder into $g$'s coordinates,
and apply $g$. Often the composite simplifies to a layout; we show that when it
does, there is a coarsest such form, whose digits every other one refines, and
we emit that.

A warp's or thread's slice is a \emph{restriction}. A restricted layout is no
longer composable, which forces the user to settle the whole view of a tile
before taking slices.

We validate the algebra by reproducing within it CuTe's operations, such as
composition, complement, and division. We implement it in OCaml and use the
implementation to rederive CuTe's atoms and usage patterns. The decision
procedure is proved correct in Rocq.
\end{abstract}
